\documentclass[11pt,a4paper]{article}
\usepackage[T1]{fontenc}
\usepackage[utf8]{inputenc}
\usepackage[margin=22mm]{geometry}
\usepackage{lmodern}
\usepackage{amsmath,amssymb}
\usepackage{enumitem}
\usepackage{xcolor}
\usepackage{microtype}
\usepackage{hyperref}
\definecolor{epflred}{HTML}{E3000B}
\hypersetup{colorlinks=true,linkcolor=epflred,
  pdftitle={ENG-654 Project 21: Reachable Workspace of Flexiv Enlight L with Fixed q3}}
\setlength{\parindent}{0pt}
\setlength{\parskip}{6pt}
\setlist[itemize]{leftmargin=*,itemsep=5pt,topsep=4pt}

\begin{document}
{\small\textbf{EPFL \textbar\ ENG-654}\hfill Kinematics-Grounded Motion Planning for Robots}
\vspace{5mm}

{\color{epflred}\Large\bfseries Project 21}

{\LARGE\bfseries Reachable Workspace of Flexiv Enlight L with Fixed $q_3$\par}

\textbf{Format:} Group of two students; 2 ECTS; simulation-based project.\\
\textbf{Prerequisites:} URDF modelling, inverse kinematics, Jacobians,
singularities, and numerical sampling.

\section*{Motivation}
Fixing the third joint of Flexiv Enlight L produces a
\textbf{6R robot parameterized by $\lambda=q_3$}, with up to eight regular
shoulder--elbow--wrist IK branches. Joint limits and singularity margins can
restrict their reachable regions differently. A workspace picture without
branch labels can hide these differences.

\section*{Task}
\textbf{Overall objective.} Compare the reachable position workspaces of all
eight branches at several fixed $q_3$ values. Hold flange orientation fixed
within each map: these are orientation-conditioned workspaces, not the complete
six-dimensional pose workspace.

\begin{itemize}
  \item \textbf{Construct and validate the 6R family.} Identify the base, joint
axes, physical limits, and flange frame in the supplied model.
Let $u=(q_1,q_2,q_4,q_5,q_6,q_7)$ and define
$F_\lambda(u)=F(q_1,q_2,\lambda,q_4,q_5,q_6,q_7)$.
Keep $\lambda$ fixed within each map. Compare reduced and original URDF
forward kinematics at 20 configurations per parameter value. Remove column 3
from the full geometric Jacobian to obtain the $6\times6$ reduced Jacobian;
check it by finite differences.

  \item \textbf{Enumerate and identify the eight branches.} Use the geometric
construction below to implement the positioning-arm and spherical-wrist IK.
Use persistent three-sign labels, retain all eight slots when solutions are
absent, and verify full-pose closure. Test three diagnostic poses and shuffled
solver output. Near a separator, flag ambiguity or singularity instead of
forcing a sign. A general-purpose 7R solver is not required.

  \item \textbf{Build comparable workspace maps.} At $\lambda=30^\circ$, choose
two distinct fixed flange orientations from validated configurations and
record their numerical rotation matrices. Evaluate the IK branches on a
common three-dimensional position grid. For each branch, distinguish absent
real IK, excessive pose residual, joint-limit violation, and insufficient
regularity. Produce eight aligned panels and a map counting how many
branches are feasible at each position. Show empty branch workspaces explicitly.

  \item \textbf{Measure the effect of fixing $q_3$.} Repeat the study at
$\lambda=-30^\circ$ and $60^\circ$ for the first orientation only, using the
same base frame, position grid, and tolerances. Compare each branch's occupied
volume, the union of all branches, and their pairwise overlap. Explain
differences using geometry, joint limits, and the branch-specific failure masks;
do not assume every branch or parameter value must give a distinct workspace.

  \item \textbf{Validate a decisive boundary.} Refine one boundary or narrow
region. Check at least 20 inside/near-boundary targets independently using
numerical IK or local continuation with $q_3$ locked. Validate any illustrated
connection at intermediate poses, checking continuity, limits, and reduced
regularity. Neighboring feasible voxels alone do not establish a feasible motion.
\end{itemize}

\newpage
\section*{Specifics and scope}
Use exactly three $q_3$ values, two orientations for the baseline and one
shared orientation for the other two values: four workspace cases in total.
Begin with approximately $30^3$ voxel centers per case in a documented box
covering the model's reach. Refine one case or a selected subregion to half
the voxel spacing. Estimate volume as the number of feasible voxels times
voxel volume; report resolution sensitivity rather than claiming exact volume.
Use identical grids for overlap comparisons. A convex hull must not be used
as a substitute for the reachable set because it can fill holes and gaps.

Use the joint-position limits of the supplied URDF consistently; distinguish
them from any separate safety limits. Resolve equivalent angles modulo $2\pi$
without discarding physically admissible representatives, and retain continuous
representatives when checking motion. Divide the translational Jacobian rows
by a stated characteristic length $\ell>0$ before comparing singular values. Record
the position, orientation, and singularity tolerances and vary the regularity
threshold in the selected refinement study.

Here, the eight requested ``aspects'' are operationally identified by the
eight fixed-$q_3$ shoulder--elbow--wrist branches. Distinguish these labels from
connected components of the limit-feasible set; filtering can disconnect a
branch. A singularity of the reduced $6\times6$ Jacobian need not be a
singularity of the full $6\times7$ Jacobian, which has no determinant.
The core study is kinematic: meshes support geometric interpretation, but
collision-free, dynamic, and hardware feasibility are outside scope.

\section*{Geometric starting point}
The robot description (.urdf) and link meshes (.stl) will be provided to the
students. Implement the analysis in a robotics library or your own code;
no additional solver or starter dataset is assumed. Verify the following
dimensions and frame conventions against the supplied model before using them.
For the flange pose $(p,R)$, define the shoulder center $S$, wrist-center
vector $w$, and arm lengths in metres:
\[
S=(0,0,0.260)^T,\qquad a=0.290,\quad b=0.280,\qquad
w=p-0.182R e_z-S,\quad e_z=(0,0,1)^T.
\]
At fixed $\lambda=q_3$, the two elbow choices and two shoulder choices follow from
\[
\begin{aligned}
c_4&=\frac{\|w\|^2-a^2-b^2}{2ab}, &q_4&=\pm\arccos(c_4),\\
v&=(b\sin q_4\cos\lambda,\ b\sin q_4\sin\lambda,\ a+b\cos q_4)^T,
&x&=\pm\sqrt{w_x^2+w_y^2-v_y^2},\\
q_1&=\operatorname{atan2}(w_y,w_x)-\operatorname{atan2}(v_y,x),\\
q_2&=\operatorname{atan2}(w_z,x)-\operatorname{atan2}(v_z,v_x).
\end{aligned}
\]
Reject non-real cases, allowing only a declared round-off tolerance. Recover
the two wrist flips from the residual orientation, accounting for the fixed
URDF frame rotations. Label a regular solution by
$(\operatorname{sgn}x,\operatorname{sgn}\sin q_4,\operatorname{sgn}\sin q_6)$.
Check the reduced determinant's proportionality to
$abx\sin q_4\sin q_6$ under your Jacobian convention. The shoulder separator
is $x=0$, not simply $q_2=0$.

\section*{Expected outcome}
Deliver reproducible code, the selected parameters and orientations, a
branch-labelled workspace atlas, volume and overlap tables, and a refinement
and boundary-validation example. Submit a concise 5--6-page report and a short
simulation demonstration. Explain which conclusions concern the fixed-$q_3$
6R family and which cannot be extended to the unrestricted seven-joint robot.

\end{document}
